% ============================================================
%  CHAPTER 11: DENSITY, PRESSURE AND ARCHIMEDES' PRINCIPLE
%  The Physics Codex: A Complete University Foundation
%  Korede Samuel Omotosho (Falcon)
%  Aurikrex Learning Systems, 2026
%
%  File: chapters/part2/ch11_density_pressure.tex
% ============================================================

\chapter{Density, Pressure and Archimedes' Principle}
\label{ch:density_pressure}
\minitoc
\newpage

\chapterquote{The difference between stupidity and genius
is that genius has its limits.}{Albert Einstein}

% ============================================================
% CHAPTER OPENING INTUITION
% ============================================================

\begin{tcolorbox}[
  enhanced,
  colback=codexblue!5,
  colframe=codexblue,
  arc=4mm,
  boxrule=1pt,
  left=10pt, right=10pt, top=10pt, bottom=10pt
]

\textit{Close your eyes for a moment and imagine two things.}

\textit{First: a small coin. You hold it over a glass of
water and let go. It drops straight through and hits the
bottom without hesitation.}

\textit{Second: an ocean liner. Half a million tonnes of
steel. Imagine placing it gently on the surface of the sea.
It sits there. Calmly. Perfectly comfortable. Going nowhere.}

\textit{The coin is made of metal. The ship is made of metal.
Both are denser than water. So why does one sink and the
other float? The answer is not in the material. It is in the
\textbf{shape}, and what that shape forces the water to do.
That insight --- that it is the water's reaction, not the
object's weight, that determines whether something sinks or
floats --- is the heart of Archimedes' principle, and it will
change how you see every floating, sinking, and submerged
object for the rest of your life.}

\textit{But before we get there, we need to understand two
simpler ideas: how tightly matter is packed (density), and
how forces spread themselves over areas (pressure). Build
those two tools, and Archimedes becomes obvious.}

\end{tcolorbox}

\vspace{0.5cm}

\begin{objectivebox}
By the end of this chapter, you should be able to:
\begin{enumerate}[noitemsep]
  \item Define density and calculate it for solids,
  liquids and gases
  \item Define pressure and apply $P = F/A$ to
  everyday situations
  \item Derive and apply the pressure-depth formula
  $P = \rho g h$
  \item State and apply Pascal's principle with
  real applications
  \item State Archimedes' principle and the law of
  floatation
  \item Calculate upthrust and determine whether
  objects float or sink
  \item Define and calculate relative density
  \item Apply these concepts to solve problems
  involving floating, submerged, and partially
  submerged objects
\end{enumerate}
\end{objectivebox}

% ============================================================
\section{Density}
% ============================================================

\subsection{The Intuition Behind Density}

Imagine you have two identical boxes. You fill one
with air and the other with sand. They look the same
from the outside. Same size. Same shape. Same volume.
But pick them up and you immediately know which is
which. The sand box is \textit{heavier} --- not because
it is bigger, but because its contents are more tightly
packed. More mass in the same space.

That ratio --- how much mass is squeezed into a given
volume --- is what we call \textbf{density}.

A piece of lead feels heavier than a piece of wood of
the same size because lead has more mass in the same
volume. It is more densely packed. A balloon filled
with air floats in your room because air inside the
balloon has roughly the same density as the air outside.
Fill the same balloon with water and it sinks --- water
is nearly 800 times denser than air.

\begin{definitionbox}[Density]
The \textbf{density} $\rho$ of a substance is the
mass per unit volume:
\[
  \rho = \frac{m}{V}
\]
SI unit: kg\,m$^{-3}$.
Also used: g\,cm$^{-3}$ (note: $1\ \text{g\,cm}^{-3}
= 1000\ \text{kg\,m}^{-3}$).

Density is a \textbf{scalar} quantity and an
\textbf{intensive} property --- it does not depend
on how much of the substance you have, only on what
it is.
\end{definitionbox}

\begin{diagbox}[Visualising Density]
\begin{center}
\begin{tikzpicture}[scale=0.85]

  % Box 1: Low density (air/wood - widely spaced dots)
  \fill[codexblue!10] (0,0) rectangle (3,3);
  \draw[thick, codexblue] (0,0) rectangle (3,3);
  \node[below, font=\small] at (1.5,-0.2) {Low density};
  % Widely spaced particles
  \foreach \x/\y in {0.5/0.5, 1.5/0.5, 2.5/0.5,
    0.5/1.5, 1.5/1.5, 2.5/1.5,
    0.5/2.5, 1.5/2.5, 2.5/2.5}
    \fill[codexblue!70] (\x,\y) circle (0.18cm);

  % Box 2: Medium density (water)
  \fill[codexaccent!10] (4,0) rectangle (7,3);
  \draw[thick, codexaccent] (4,0) rectangle (7,3);
  \node[below, font=\small] at (5.5,-0.2) {Medium density};
  % Medium spaced particles
  \foreach \x/\y in {
    4.35/0.35, 4.85/0.35, 5.35/0.35, 5.85/0.35, 6.35/0.35,
    4.35/0.85, 4.85/0.85, 5.35/0.85, 5.85/0.85, 6.35/0.85,
    4.35/1.35, 4.85/1.35, 5.35/1.35, 5.85/1.35, 6.35/1.35,
    4.35/1.85, 4.85/1.85, 5.35/1.85, 5.85/1.85, 6.35/1.85,
    4.35/2.35, 4.85/2.35, 5.35/2.35, 5.85/2.35, 6.35/2.35,
    4.35/2.75, 4.85/2.75, 5.35/2.75, 5.85/2.75, 6.35/2.75}
    \fill[codexaccent!70] (\x,\y) circle (0.12cm);

  % Box 3: High density (steel/lead)
  \fill[codexred!10] (8,0) rectangle (11,3);
  \draw[thick, codexred] (8,0) rectangle (11,3);
  \node[below, font=\small] at (9.5,-0.2) {High density};
  % Very tightly packed particles
  \foreach \x in {8.2, 8.53, 8.86, 9.19, 9.52, 9.85, 10.18, 10.51, 10.84}
    \foreach \y in {0.2, 0.53, 0.86, 1.19, 1.52, 1.85, 2.18, 2.51, 2.84}
      \fill[codexred!70] (\x,\y) circle (0.13cm);

  % Same volume label
  \draw[<->, gray, thick] (0,3.4) -- (3,3.4);
  \draw[<->, gray, thick] (4,3.4) -- (7,3.4);
  \draw[<->, gray, thick] (8,3.4) -- (11,3.4);
  \node[above, gray, font=\small] at (1.5,3.4)
    {Same volume};
  \node[above, gray, font=\small] at (5.5,3.4)
    {Same volume};
  \node[above, gray, font=\small] at (9.5,3.4)
    {Same volume};

\end{tikzpicture}
\end{center}
\small\textit{All three boxes have the same volume.
More particles packed into the same space = higher
density = more mass.}
\end{diagbox}

\begin{table}[H]
\centering
\caption{Densities of Common Substances}
\label{tab:densities}
\begin{tabular}{lll}
\toprule
\textbf{Substance} & \textbf{Density (kg\,m$^{-3}$)} &
\textbf{Notes}\\
\midrule
Air (at sea level, 20°C) & 1.20 & Varies with altitude\\
Cork & 200--300 & Floats in water\\
Wood (typical) & 400--900 & Most wood floats\\
Ice & 917 & Floats in water (just)\\
Water (at 4°C) & 1000 & Reference standard\\
Sea water & 1020--1030 & Saltier = denser\\
Concrete & 2300 & Sinks in water\\
Aluminium & 2700 & Light metal\\
Steel & 7800 & \\
Copper & 8900 & \\
Lead & 11\,340 & Very dense\\
Gold & 19\,300 & Dense enough to stay in rivers\\
\bottomrule
\end{tabular}
\end{table}

\begin{examplebox}[11.1]
\stepcounter{workedexample}
A rectangular block of aluminium measures
$10\ \text{cm} \times 6\ \text{cm} \times 4\ \text{cm}$
and has mass $648\ \text{g}$.\\
(a) Calculate the density in kg\,m$^{-3}$.\\
(b) Identify the material and verify against the table.
\end{examplebox}

\begin{solutionbox}
\textbf{(a)}
\[
  V = 0.10 \times 0.06 \times 0.04
  = 2.4\times10^{-4}\ \text{m}^3
\]
\[
  \rho = \frac{m}{V}
  = \frac{0.648}{2.4\times10^{-4}}
  = 2700\ \text{kg\,m}^{-3}
\]

\textbf{(b)} From Table~\ref{tab:densities},
$\rho = 2700\ \text{kg\,m}^{-3}$ matches
\textbf{aluminium} exactly.
\end{solutionbox}

% ============================================================
\section{Pressure}
% ============================================================

\subsection{The Intuition Behind Pressure}

Think about this carefully. You are standing in soft
mud. You are wearing flat sandals and you barely sink.
Now you hand your sandals to a friend and balance on
the tips of two pencils instead. You suddenly sink
deeply into the mud. Your weight has not changed.
Your mass has not changed. The force pushing down on
the mud is identical in both cases.

What changed is the \textit{area} over which that
force is spread.

The flat sandal spreads your weight over a large area
--- the mud receives a small force per square centimetre
and holds. The pencil tip concentrates your entire
weight onto a tiny area --- the mud receives an enormous
force per square centimetre and gives way. That force
per unit area is \textbf{pressure}.

This is why a nurse uses a needle (tiny area, high
pressure) to pierce skin but a doctor uses a flat
bandage (large area, low pressure) to stop bleeding.
Same idea, opposite purposes.

\begin{definitionbox}[Pressure]
\textbf{Pressure} is the force per unit area acting
perpendicular to a surface:
\[
  P = \frac{F}{A}
\]
SI unit: N\,m$^{-2}$ = Pascal (Pa).

Pressure is a \textbf{scalar} quantity. The same force
produces higher pressure over a smaller area and lower
pressure over a larger area.
\end{definitionbox}

\begin{diagbox}[Pressure Depends on Area]
\begin{center}
\begin{tikzpicture}[scale=0.9]

  % Ground (soft surface)
  \fill[brown!30] (0,-0.5) rectangle (12,0);
  \draw[thick, brown!60] (0,0) -- (12,0);
  % Deformation lines
  \draw[brown!50, thick] (1.5,0) -- (1.5,-0.35);
  \draw[brown!50, thick] (2.5,0) -- (2.5,-0.35);
  \draw[brown!50, thick] (8.5,0) -- (8.7,-0.5);
  \draw[brown!50, thick] (9.5,0) -- (9.3,-0.5);

  % FLAT SHOE (left)
  % Shoe sole (wide area)
  \fill[codexblue!40] (0.8,0) rectangle (3.2,0.3);
  \draw[thick, codexblue] (0.8,0) rectangle (3.2,0.3);
  % Leg
  \draw[thick] (2.0,0.3) -- (2.0,2.0);
  % Person
  \fill[codexblue!30] (1.5,2.0) rectangle (2.5,3.0);
  \node at (2.0,2.5) {$F$};
  % Force arrow down
  \draw[->, very thick, codexblue]
    (2.0,2.0) -- (2.0,0.3);
  % Area label
  \draw[<->, codexblue, thick] (0.8,-0.8) -- (3.2,-0.8);
  \node[below, codexblue] at (2.0,-0.9)
    {Large area $A_1$};
  % Pressure label
  \node[codexblue, font=\small] at (2.0,-2.0)
    {Small $P = F/A_1$};
  \node[below, font=\small, gray] at (2.0,-2.5)
    {Barely sinks};

  % PENCIL (right, same force)
  % Pencil tip (tiny area)
  \fill[codexred!40] (8.7,0) rectangle (9.3,0.3);
  \draw[thick, codexred] (8.7,0) rectangle (9.3,0.3);
  % Leg
  \draw[thick] (9.0,0.3) -- (9.0,2.0);
  % Person (same force)
  \fill[codexred!30] (8.5,2.0) rectangle (9.5,3.0);
  \node at (9.0,2.5) {$F$};
  % Force arrow
  \draw[->, very thick, codexred]
    (9.0,2.0) -- (9.0,0.3);
  % Area label
  \draw[<->, codexred, thick] (8.7,-0.8) -- (9.3,-0.8);
  \node[below, codexred] at (9.0,-0.9)
    {Tiny area $A_2$};
  % Pressure label
  \node[codexred, font=\small] at (9.0,-2.0)
    {Huge $P = F/A_2$};
  \node[below, font=\small, gray] at (9.0,-2.5)
    {Sinks deeply};

  % Same force label
  \node[font=\small, align=center] at (5.5,2.5)
    {Same force $F$\\$\downarrow$};

\end{tikzpicture}
\end{center}
\small\textit{Identical forces, completely different
pressures. The area is everything.}
\end{diagbox}

\begin{examplebox}[11.2]
\stepcounter{workedexample}
A $70\ \text{kg}$ person stands on one foot. The area
of the foot in contact with the ground is
$150\ \text{cm}^2 = 0.015\ \text{m}^2$.\\
(a) Calculate the pressure on the ground.\\
(b) The same person balances on the tip of a high-heel
shoe of area $1.0\ \text{cm}^2 = 1.0\times10^{-4}\ \text{m}^2$.
Calculate the new pressure.\\
(c) Compare with atmospheric pressure
($\approx 1.01\times10^5\ \text{Pa}$).
($g = 10\ \text{m\,s}^{-2}$)
\end{examplebox}

\begin{solutionbox}
Weight: $F = mg = 70 \times 10 = 700\ \text{N}$

\textbf{(a)} Flat foot:
\[
  P = \frac{F}{A} = \frac{700}{0.015}
  = 46\,667\ \text{Pa} \approx 47\ \text{kPa}
\]

\textbf{(b)} High heel:
\[
  P = \frac{700}{1.0\times10^{-4}}
  = 7\times10^6\ \text{Pa} = 7\ \text{MPa}
\]

\textbf{(c)} The heel pressure is
$7\times10^6 / 1.01\times10^5 \approx 69$ times
atmospheric pressure. The pressure from a high-heel
shoe is equivalent to roughly 70 atmospheres --- more
than enough to damage soft flooring permanently.
This is why wooden floors and aircraft carpets have
signs prohibiting stiletto heels.
\end{solutionbox}

% ============================================================
\section{Pressure in Fluids}
% ============================================================

\subsection{The Intuition: Why Pressure Increases with Depth}

Imagine you are swimming underwater. You dive
shallowly --- you feel fine. You dive deeper ---
you feel your ears start to hurt. You go deeper
still and the pressure becomes painful. Why?

Think about what is above you at each depth.
When you are just below the surface, there is
only a thin layer of water pushing down on you.
When you are $10\ \text{m}$ deep, there is a
$10\ \text{m}$-tall column of water above you,
plus all the atmosphere on top of that. The
deeper you go, the more water is stacked above
you, the more weight is pressing down, and the
higher the pressure.

This is not just true in the downward direction.
Fluids transmit pressure in \textit{all directions}
simultaneously --- which is what makes fluid
pressure unique compared to the pressure from
a solid object resting on a surface.

\begin{processbox}[Deriving Pressure at Depth $h$]
Consider a horizontal surface of area $A$ at depth $h$
below the surface of a fluid of density $\rho$.

The volume of fluid above this surface:
$V = Ah$

The mass of this fluid column:
$m = \rho V = \rho A h$

The weight of this fluid column (force on the surface):
$F = mg = \rho A h g$

The pressure due to this column:
\[
  P = \frac{F}{A} = \frac{\rho A h g}{A}
  = \boxed{\rho g h}
\]

Adding atmospheric pressure $P_0$ for the absolute
(total) pressure:
\[
  P_{total} = P_0 + \rho g h
\]
\end{processbox}

\begin{diagbox}[Pressure Increases with Depth]
\begin{center}
\begin{tikzpicture}[scale=0.9]
  % Water container
  \fill[codexaccent!15] (0,0) rectangle (8,6);
  \draw[very thick, codexblue!60] (0,0) rectangle (8,6);
  % Surface label
  \draw[thick, codexaccent] (0,6) -- (8,6);
  \node[above, codexaccent] at (4,6.1)
    {Water surface ($P = P_0$)};

  % Depth brackets (drawn longest-to-shortest so each color stays visible
  % in the band where its own label sits)
  \draw[<->, codexred, thick] (-0.5,1) -- (-0.5,6);
  \node[left, codexred, font=\small] at (-0.6,2.6) {$h_3$};
  \draw[<->, codexgreen, thick] (-0.5,3) -- (-0.5,6);
  \node[left, codexgreen, font=\small] at (-0.6,4.1) {$h_2$};
  \draw[<->, codexblue, thick] (-0.5,5) -- (-0.5,6);
  \node[left, codexblue, font=\small] at (-0.6,5.8) {$h_1$};

  % Depth 1 (shallow) - thin arrows, fixed length 0.6
  \draw[->, line width=1pt, codexblue] (1,5) -- (1.6,5);
  \draw[->, line width=1pt, codexblue] (1,5) -- (0.4,5);
  \draw[->, line width=1pt, codexblue] (1,5) -- (1,5.6);
  \draw[->, line width=1pt, codexblue] (1,5) -- (1,4.4);
  \node[right, font=\tiny, codexblue] at (1.7,5)
    {$P_1 = P_0 + \rho g h_1$};

  % Depth 2 (medium) - thicker arrows, same length 0.6
  \draw[->, line width=1.6pt, codexgreen] (1,3) -- (1.6,3);
  \draw[->, line width=1.6pt, codexgreen] (1,3) -- (0.4,3);
  \draw[->, line width=1.6pt, codexgreen] (1,3) -- (1,3.6);
  \draw[->, line width=1.6pt, codexgreen] (1,3) -- (1,2.4);
  \node[right, font=\tiny, codexgreen] at (1.7,3)
    {$P_2 = P_0 + \rho g h_2$};

  % Depth 3 (deep) - thickest arrows, same length 0.6
  \draw[->, line width=2.4pt, codexred] (1,1) -- (1.6,1);
  \draw[->, line width=2.4pt, codexred] (1,1) -- (0.4,1);
  \draw[->, line width=2.4pt, codexred] (1,1) -- (1,1.6);
  \draw[->, line width=2.4pt, codexred] (1,1) -- (1,0.4);
  \node[right, font=\tiny, codexred] at (1.7,1)
    {$P_3 = P_0 + \rho g h_3$};

  % Fluid column visualisation
  \fill[codexblue!30, opacity=0.6]
    (6.5,1) rectangle (7.5,6);
  \draw[thick, dashed, codexblue]
    (6.5,1) rectangle (7.5,6);
  \draw[->, thick] (7.0,5.5) -- (7.0,1.2);
  \node[right, font=\tiny] at (7.5,3.5)
    {\shortstack[l]{Column\\of fluid\\above}};
\end{tikzpicture}
\end{center}
\small\textit{The deeper the point, the taller the
column of fluid above it, the greater the pressure.
Pressure also acts equally in \textbf{all directions}
at any given depth --- that's why the arrows at each
depth point outward on all four sides.}
\end{diagbox}

\begin{keybox}[Key Properties of Fluid Pressure]
\begin{itemize}[noitemsep]
  \item Pressure in a fluid acts equally in
  \textbf{all directions} at any point
  \item Pressure depends only on \textbf{depth},
  \textbf{density} and $g$ --- not on the shape
  of the container, not on the total volume of fluid
  \item At the same depth in a connected fluid,
  pressure is the same --- even if the containers
  are different shapes (this is the paradox of
  hydrostatics)
  \item Pressure \textbf{increases} linearly with depth:
  $P = P_0 + \rho g h$
\end{itemize}
\end{keybox}

\begin{diagbox}[The Hydrostatic Paradox: Same Depth = Same Pressure]
\begin{center}
\begin{tikzpicture}[scale=0.85]

  % Separate vessels: each has the same liquid and free-surface height.
  % The equal-depth comparison is measured below those free surfaces.
  \fill[codexaccent!20] (0.5,0.5) rectangle (2.0,4.5);
  \draw[thick] (0.5,0.5) -- (0.5,4.5);
  \draw[thick] (2.0,0.5) -- (2.0,4.5);
  \draw[thick, codexaccent] (0.5,4.5) -- (2.0,4.5);
  \draw[thick] (0.5,0.5) -- (2.0,0.5);
  \node[above, font=\small] at (1.25,4.6) {Narrow};

  \fill[codexaccent!20] (3.0,0.5) rectangle (6.5,4.5);
  \draw[thick] (3.0,0.5) -- (3.0,4.5);
  \draw[thick] (6.5,0.5) -- (6.5,4.5);
  \draw[thick, codexaccent] (3.0,4.5) -- (6.5,4.5);
  \draw[thick] (3.0,0.5) -- (6.5,0.5);
  \node[above, font=\small] at (4.75,4.6) {Wide};

  \fill[codexaccent!20]
    (8.0,0.5) -- (7.5,4.5) -- (10.0,4.5) -- (9.5,0.5) -- cycle;
  \draw[thick] (8.0,0.5) -- (7.5,4.5);
  \draw[thick] (9.5,0.5) -- (10.0,4.5);
  \draw[thick, codexaccent] (7.5,4.5) -- (10.0,4.5);
  \draw[thick] (8.0,0.5) -- (9.5,0.5);
  \node[above, font=\small] at (8.75,4.6) {Flared};

  % Same depth h below each free surface (y=4.5 to y=2.5)
  \draw[dashed, codexred, thick] (0.2,2.5) -- (10.0,2.5);
  \node[right, codexred, font=\small] at (10.1,2.5)
    {Same depth};
  \draw[<->, gray, thick] (2.35,2.5) -- (2.35,4.5);
  \node[right, gray, font=\small] at (2.4,3.5) {$h$};

  % Pressure markers at the same depth in each separate vessel
  \draw[->, codexred] (1.25,2.5) -- (0.65,2.5);
  \draw[->, codexred] (1.25,2.5) -- (1.85,2.5);
  \draw[->, codexred] (4.75,2.5) -- (4.15,2.5);
  \draw[->, codexred] (4.75,2.5) -- (5.35,2.5);
  \draw[->, codexred] (8.75,2.5) -- (8.15,2.5);
  \draw[->, codexred] (8.75,2.5) -- (9.35,2.5);

  \node[below, font=\small, align=center, text width=10cm] at (5,-0.3)
    {At depth $h$ below each free surface, pressure is
    $P = P_0 + \rho g h$ in all three vessels, regardless of shape.};

\end{tikzpicture}
\end{center}
\end{diagbox}

\begin{examplebox}[11.3]
\stepcounter{workedexample}
A diver descends to $30\ \text{m}$ below the sea
surface. The density of sea water is
$1025\ \text{kg\,m}^{-3}$ and atmospheric pressure
is $1.01\times10^5\ \text{Pa}$.
($g = 10\ \text{m\,s}^{-2}$)\\
(a) Calculate the gauge pressure at $30\ \text{m}$.\\
(b) Calculate the total (absolute) pressure.\\
(c) Express the total pressure in atmospheres.
\end{examplebox}

\begin{solutionbox}
\textbf{(a)} Gauge pressure (due to water column only):
\[
  P_{gauge} = \rho g h
  = 1025 \times 10 \times 30
  = 307\,500\ \text{Pa} \approx 3.08\times10^5\ \text{Pa}
\]

\textbf{(b)} Total pressure:
\[
  P_{total} = P_0 + \rho g h
  = 1.01\times10^5 + 3.075\times10^5
  = 4.085\times10^5\ \text{Pa}
\]

\textbf{(c)} In atmospheres:
\[
  \frac{P_{total}}{P_0}
  = \frac{4.085\times10^5}{1.01\times10^5}
  \approx 4.05\ \text{atm}
\]
Every $10\ \text{m}$ of sea water adds approximately
1 atmosphere of pressure. At $30\ \text{m}$, the
total pressure is about 4 atmospheres --- which is
why scuba divers must use pressurised air cylinders.
\end{solutionbox}

% ============================================================
\section{Pascal's Principle}
% ============================================================

\subsection{The Intuition: Pressure as a Messenger}

Imagine squeezing a tube of toothpaste at one end.
Toothpaste comes out the other end. You applied
force at one point, and the effect appeared somewhere
else entirely. The tube transmitted your force through
the paste.

Fluids do the same thing --- but better. Pascal
discovered that when you apply pressure to an enclosed
fluid, the fluid does not just transmit that pressure
in one direction. It transmits it \textit{equally
in every direction throughout the entire fluid}.
The fluid acts as a perfect pressure messenger.

This is the principle behind every hydraulic system:
the car lift in a mechanic's workshop, the braking
system in every vehicle on the road, the landing gear
of aircraft, and the hydraulic press in a factory.

\begin{lawbox}[Pascal's Principle]
Pressure applied to an enclosed, incompressible fluid
is transmitted equally and undiminished to every part
of the fluid and to the walls of its container.

Mathematically: if pressure $P$ is applied at one
point in an enclosed fluid:
\[
  P_{transmitted} = P \quad
  \text{(at every other point in the fluid)}
\]

For a hydraulic system with two pistons of areas
$A_1$ (small) and $A_2$ (large):
\[
  \frac{F_1}{A_1} = \frac{F_2}{A_2}
  \implies F_2 = F_1 \times \frac{A_2}{A_1}
\]
A small force on the small piston produces a large
force on the large piston.
\end{lawbox}

\begin{diagbox}[Hydraulic Lift: Pascal's Principle in Action]
\begin{center}
\begin{tikzpicture}[scale=0.9]

  % Fluid reservoir
  \fill[codexaccent!20]
    (0,0) rectangle (10,1.5);
  \fill[codexaccent!20]
    (0,1.5) rectangle (2.5,3.5);
  \fill[codexaccent!20]
    (7,1.5) rectangle (10,5.5);
  \draw[thick, codexblue!60]
    (0,0) rectangle (10,1.5);
  \draw[thick, codexblue!60] (0,1.5) -- (2.5,1.5);
  \draw[thick, codexblue!60] (2.5,1.5) -- (2.5,4.5);
  \draw[thick, codexblue!60] (7,1.5) -- (10,1.5);
  \draw[thick, codexblue!60] (7,1.5) -- (7,7.5);

  % Small piston (left)
  \fill[gray!50] (0.5,3.3) rectangle (2.0,3.7);
  \draw[thick] (0.5,3.3) rectangle (2.0,3.7);
  \node[left, font=\small, align=left] at (0.4,3.5)
    {Small piston\\$A_1$};

  % Effort force on small piston
  \draw[->, very thick, codexblue]
    (1.25,4.5) -- (1.25,3.9)
    node[midway, right, font=\small]{$F_1$ (small)};

  % Large piston (right)
  \fill[gray!50] (7.2,5.3) rectangle (9.8,5.7);
  \draw[thick] (7.2,5.3) rectangle (9.8,5.7);
  \node[right, font=\small] at (10.1,5.5)
    [right, font=\small, align=center] at (10.1,5.5) {Large piston\\$A_2 \gg A_1$};

  % Load on large piston
  \fill[codexred!30] (7.2,5.7) rectangle (9.8,7.0);
  \draw[thick] (7.2,5.7) rectangle (9.8,7.0);
  \node at (8.5,6.3) {LOAD};
  \node[font=\small] at (8.5,7.6)
    {$\Delta F_2 = \Delta P\,A_2 = F_1(A_2/A_1)$};
  \node[font=\small] at (8.5,8.2) {(Large force!)};

  % Pressure arrows in fluid
  \foreach \x/\y in {3/0.75, 5/0.75, 7.5/0.75,
    1.25/2.0, 1.25/2.8}
    \fill[codexaccent] (\x,\y) circle (1.5pt);

  % Equal applied pressure increment label
  \node[codexaccent, font=\small] at (5,0.3)
    {Applied pressure increment $\Delta P$ is transmitted};

  % Area comparison
  \draw[<->, gray] (0.5,4.8) -- (2.0,4.8);
  \node[above, gray, font=\small] at (1.25,4.8)
    {$A_1$};
  \draw[<->, gray] (7.2,4.8) -- (9.8,4.8);
  \node[above, gray, font=\small] at (8.5,4.8)
    {$A_2$};

\end{tikzpicture}
\end{center}
\small\textit{A small force $F_1$ applied to the small
piston produces a pressure increment $\Delta P = F_1/A_1$,
which is transmitted through the fluid by Pascal's principle.
The ideal force increment on the large piston is
$\Delta F_2 = \Delta P A_2 = F_1(A_2/A_1) \gg F_1$.
Absolute pressure also varies with height in a static fluid.}
\end{diagbox}

\begin{examplebox}[11.4]
\stepcounter{workedexample}
A hydraulic car lift has a small piston of diameter
$2.0\ \text{cm}$ and a large piston of diameter
$20\ \text{cm}$. A force of $100\ \text{N}$ is
applied to the small piston.\\
(a) Calculate the force exerted by the large piston.\\
(b) If a $1500\ \text{kg}$ car is to be lifted,
what force must be applied to the small piston?
($g = 10\ \text{m\,s}^{-2}$)
\end{examplebox}

\begin{solutionbox}
\[
  A_1 = \pi(0.01)^2 = 3.14\times10^{-4}\ \text{m}^2
\]
\[
  A_2 = \pi(0.10)^2 = 3.14\times10^{-2}\ \text{m}^2
\]

\textbf{(a)}
\[
  F_2 = F_1 \times \frac{A_2}{A_1}
  = 100 \times \frac{3.14\times10^{-2}}
  {3.14\times10^{-4}}
  = 100 \times 100 = 10\,000\ \text{N}
\]

\textbf{(b)} Weight of car: $W = 1500 \times 10
= 15\,000\ \text{N}$
\[
  F_1 = F_2 \times \frac{A_1}{A_2}
  = 15\,000 \times \frac{1}{100}
  = 150\ \text{N}
\]
A $150\ \text{N}$ effort (roughly the weight of
a $15\ \text{kg}$ mass) lifts a $1500\ \text{kg}$ car.
\end{solutionbox}

% ============================================================
\section{Archimedes' Principle}
% ============================================================

\subsection{The Intuition: The Displaced Water Fights Back}

Here is the core idea. When you push an object into
water, you are forcing the water out of the way.
The water has to go somewhere --- it moves aside
and upward. But water does not move without
resistance. The water that has been displaced
\textit{pushes back} on whatever displaced it.
That pushback is the upthrust.

Now here is the beautiful part. How hard does the
water push back? Exactly as hard as the weight of
the water that was displaced. Not more, not less.
Precisely the weight of the water that the object
forced out of the way.

A small object displaces a small amount of water
and receives a small upthrust. A large object
displaces a large amount of water and receives
a large upthrust. The object floats if its weight
is less than or equal to the maximum upthrust
(the weight of fluid it would displace when fully
submerged). It sinks if its weight exceeds the
maximum upthrust.

\begin{lawbox}[Archimedes' Principle]
When a body is wholly or partially immersed in a
fluid, it experiences an \textbf{upthrust} (buoyant
force) equal to the weight of the fluid displaced:
\[
  U = \rho_{fluid} \times V_{displaced} \times g
  = W_{fluid\ displaced}
\]
The upthrust acts vertically upward through the
\textbf{centre of buoyancy} (the centroid of the
displaced volume).
\end{lawbox}

\begin{diagbox}{Archimedes' Principle --- Step by Step}
\begin{center}
\begin{tikzpicture}[scale=0.8]
  % ---------- STEP 1: Before immersion ----------
  \node[font=\bfseries\small, codexblue, align=center, text width=3.6cm]
    at (2.25,9) {Step 1: Before immersion};
  \draw[thick] (0,0) rectangle (4.5,6);
  \fill[codexaccent!20] (0,0) rectangle (4.5,3.5);
  \draw[thick, codexaccent] (0,3.5) -- (4.5,3.5);
  \node[anchor=west, font=\small] at (0.15,3.65) {Water level};
  \fill[codexblue!40] (1.45,6.6) rectangle (3.05,8.2);
  \draw[thick] (1.45,6.6) rectangle (3.05,8.2);
  \node at (2.25,7.4) {Object};
  \draw[->, very thick, codexred]
    (2.25,6.6) -- (2.25,5.9)
    node[midway, right, font=\small] {$W = mg$};
  % ---------- STEP 2: Partial immersion ----------
  \node[font=\bfseries\small, codexblue, align=center, text width=3.6cm]
    at (8.25,9) {Step 2: Partial immersion};
  \draw[thick] (6,0) rectangle (10.5,6);
  \draw[thick, dashed, gray] (6,3.5) -- (10.5,3.5);
  \fill[codexaccent!20] (6,0) rectangle (10.5,3.9);
  \draw[thick, codexaccent] (6,3.9) -- (10.5,3.9);
  \node[anchor=west, font=\small] at (6.15,5.0) {Water rises};
  \fill[codexblue!40] (7.45,3.9) rectangle (9.05,4.7);
  \fill[codexblue!70, opacity=0.5] (7.45,3.1) rectangle (9.05,3.9);
  \draw[thick] (7.45,3.1) rectangle (9.05,4.7);
  \node at (8.25,4.3) {Object};
  \draw[->, very thick, codexred]
    (9.4,4.7) -- (9.4,3.9)
    node[midway, right, font=\small] {$W$};
  \draw[->, very thick, codexgreen]
    (7.1,3.1) -- (7.1,3.9)
    node[midway, left, font=\small] {$U$};
  % ---------- STEP 3: Full immersion ----------
  \node[font=\bfseries\small, codexblue, align=center, text width=3.6cm]
    at (14.25,9) {Step 3: Full immersion};
  \draw[thick] (12,0) rectangle (16.5,6);
  \draw[thick, dashed, gray] (12,3.5) -- (16.5,3.5);
  \fill[codexaccent!20] (12,0) rectangle (16.5,4.3);
  \draw[thick, codexaccent] (12,4.3) -- (16.5,4.3);
  \node[anchor=east, font=\small] at (16.35,4.45) {Water displaced};
  \fill[codexblue!40] (13.45,2.5) rectangle (15.05,4.1);
  \draw[thick] (13.45,2.5) rectangle (15.05,4.1);
  \node at (14.25,3.3) {Object};
  \draw[->, codexaccent, thick] (13.1,3.3) -- (13.45,3.3);
  \draw[->, codexaccent, thick] (15.4,3.3) -- (15.05,3.3);
  \draw[->, very thick, codexred]
    (15.4,4.1) -- (15.4,2.8)
    node[midway, right, font=\small] {$W$};
  \draw[->, very thick, codexgreen]
    (13.1,2.5) -- (13.1,3.9)
    node[midway, left, font=\small] {$U$};
  \node[font=\small, align=center, codexgreen]
    at (14.25,1.4)
    {$U = \rho_{fluid}\,V_{obj}\,g$\\(weight of displaced fluid)};
\end{tikzpicture}
\end{center}
\end{diagbox}

\begin{lawbox}[Law of Floatation]
A floating body displaces its own \textbf{weight}
of the fluid in which it floats:
\[
  W_{object} = U = \rho_{fluid} \times V_{submerged}
  \times g
\]
For a floating object:
\begin{itemize}[noitemsep]
  \item If $\rho_{object} < \rho_{fluid}$: object
  floats (upthrust when fully submerged exceeds weight)
  \item If $\rho_{object} = \rho_{fluid}$: object
  is neutrally buoyant (floats submerged at any depth)
  \item If $\rho_{object} > \rho_{fluid}$: object
  sinks (weight always exceeds maximum upthrust)
\end{itemize}
\end{lawbox}

\begin{diagbox}[Floating, Sinking and Neutral Buoyancy]
\begin{center}
\begin{tikzpicture}[scale=0.85]

  % Water in container
  \fill[codexaccent!15] (0,0) rectangle (12,5);
  \draw[thick] (0,0) rectangle (12,5);
  \draw[thick, codexaccent] (0,5) -- (12,5);

  % FLOATS (left): wood, $\rho < \rho_{water}$
  \fill[codexblue!60] (1,4.5) rectangle (2.8,5.5);
  \draw[thick] (1,4.5) rectangle (2.8,5.5);
  \node at (1.9,5.0) {\small Wood};
  % Only partially submerged
  \fill[codexblue!80, opacity=0.5]
    (1,4.5) rectangle (2.8,5.0);
  \draw[->, codexgreen, thick]
    (1.3,4.55) -- (1.3,5.45)
    node[above, font=\tiny]{$U=W$};
  \draw[->, codexred, thick]
    (2.5,5.45) -- (2.5,4.55)
    node[below, font=\tiny]{$W=U$};
  \node[below, font=\small] at (1.9,2.3)
    {Floats};
  \node[below, font=\tiny, codexblue] at (1.9,3.0)
    {$\rho < \rho_{fluid}$};

  % NEUTRAL (middle): same density
  \fill[codexaccent!60] (5,2.3) rectangle (7,3.5);
  \draw[thick] (5,2.3) rectangle (7,3.5);
  \node at (6,2.9) {\small Submarine};
  \draw[->, codexgreen, thick]
    (5.4,2.35) -- (5.4,3.25)
    node[above, font=\tiny]{$U=W$};
  \draw[->, codexred, thick]
    (6.6,3.45) -- (6.6,2.55)
    node[below, font=\tiny]{$W=U$};
  \node[below, font=\small] at (6,0.1)
    {Neutrally buoyant};
  \node[below, font=\tiny, codexaccent] at (6,0.8)
    {$\rho = \rho_{fluid}$};

  % SINKS (right): metal
  \fill[codexred!50] (9.5,0.2) rectangle (11,1.2);
  \draw[thick] (9.5,0.2) rectangle (11,1.2);
  \node at (10.25,0.7) {\small Metal};
  \draw[->, codexgreen, thick]
    (9.65,0.2) -- (9.65,0.9)
    node[above, font=\tiny]{$U$};
  \draw[->, codexred, thick, line width=2pt]
    (10.85,1.2) -- (10.85,0.2)
    node[below, font=\tiny]{$W > U$};
  \node[below, font=\small] at (10.25,-0.8)
    {Sinks};
  \node[below, font=\tiny, codexred] at (10.25,3.4)
    {$\rho > \rho_{fluid}$};

\end{tikzpicture}
\end{center}
\end{diagbox}

\begin{examplebox}[11.5]
\stepcounter{workedexample}
A metal crown of mass $420\ \text{g}$ appears to have
mass $370\ \text{g}$ when weighed while fully submerged
in water ($\rho_w = 1000\ \text{kg\,m}^{-3}$).\\
(a) Calculate the upthrust.\\
(b) Calculate the volume of the crown.\\
(c) Calculate the density of the crown material.\\
(d) Is the crown made of gold ($\rho_{gold}
= 19\,300\ \text{kg\,m}^{-3}$)?
($g = 10\ \text{m\,s}^{-2}$)
\end{examplebox}

\begin{solutionbox}
\textbf{(a)} Upthrust = loss in apparent weight:
\[
  U = W_{air} - W_{water}
  = (420 - 370) \times 10^{-3} \times 10
  = 0.50\ \text{N}
\]

\textbf{(b)} From $U = \rho_w V g$:
\[
  V = \frac{U}{\rho_w g}
  = \frac{0.50}{1000 \times 10}
  = 5.0\times10^{-5}\ \text{m}^3
\]

\textbf{(c)} Density:
\[
  \rho_{crown} = \frac{m}{V}
  = \frac{420\times10^{-3}}{5.0\times10^{-5}}
  = 8400\ \text{kg\,m}^{-3}
\]

\textbf{(d)} The density $8400\ \text{kg\,m}^{-3}$
is close to copper or bronze --- far less than gold's
$19\,300\ \text{kg\,m}^{-3}$. \textbf{The crown is
not gold.} This is precisely the problem Archimedes
was asked to solve for King Hiero II of Syracuse ---
and the principle he used to solve it.
\end{solutionbox}

% ============================================================
\section{Relative Density}
% ============================================================

\subsection{The Intuition: Comparing to Water}

Saying a material has density $7800\ \text{kg\,m}^{-3}$
is informative if you know what $\text{kg\,m}^{-3}$
means physically. But saying it is 7.8 times denser
than water is immediately intuitive --- everyone knows
how heavy water feels. Relative density is that
comparison: density of the substance expressed as a
multiple of water's density.

\begin{definitionbox}[Relative Density]
The \textbf{relative density} (RD) of a substance
is the ratio of its density to the density of water:
\[
  RD = \frac{\rho_{substance}}{\rho_{water}}
  = \frac{\rho_{substance}}{1000\ \text{kg\,m}^{-3}}
\]
RD is \textbf{dimensionless} (no unit).
It equals the density in g\,cm$^{-3}$.

RD can also be measured experimentally:
\[
  RD = \frac{W_{air}}{W_{air} - W_{water}}
  = \frac{\text{weight in air}}
  {\text{upthrust in water}}
\]
\end{definitionbox}

% ============================================================
\section{Chapter Summary}
% ============================================================

\begin{summarybox}
\textbf{Density:} $\rho = m/V$ (kg\,m$^{-3}$)

\textbf{Pressure:} $P = F/A$ (Pa = N\,m$^{-2}$)

\textbf{Pressure at depth:}
$P = P_0 + \rho g h$ (increases linearly with depth)

\textbf{Pascal's principle:} Pressure applied to an
enclosed fluid is transmitted equally in all directions.
$F_2/F_1 = A_2/A_1$

\textbf{Archimedes' principle:}
$U = \rho_{fluid}\,V_{displaced}\,g$

\textbf{Floatation condition:}
$W_{object} = U$ (weight of fluid displaced)

Floats if $\rho_{obj} < \rho_{fluid}$;
sinks if $\rho_{obj} > \rho_{fluid}$.

\textbf{Relative density:}
$RD = \rho_{substance}/\rho_{water}$;
$RD = W_{air}/(W_{air}-W_{water})$
\end{summarybox}

% ============================================================
\newpage
\begin{exercisebox}[11]

\subsection*{Section A: Multiple Choice}

\textbf{A1.} A force of $60\ \text{N}$ acts on an
area of $0.3\ \text{m}^2$. The pressure is:

\mcoption{A}{$18\ \text{Pa}$}
\mcoption{B}{$200\ \text{Pa}$}
\mcoption{C}{$0.005\ \text{Pa}$}
\mcoption{D}{$180\ \text{Pa}$}
\ansref{Ch11-A1}

\vspace{0.4cm}

\textbf{A2.} A submarine descends to a depth where
the gauge pressure is $5\times10^5\ \text{Pa}$.
The depth is (seawater density $= 1000\ \text{kg\,m}^{-3}$,
$g = 10\ \text{m\,s}^{-2}$):

\mcoption{A}{$5\ \text{m}$}
\mcoption{B}{$50\ \text{m}$}
\mcoption{C}{$500\ \text{m}$}
\mcoption{D}{$5000\ \text{m}$}
\ansref{Ch11-A2}

\vspace{0.4cm}

\textbf{A3.} An object weighs $12\ \text{N}$ in air
and $8\ \text{N}$ when fully submerged in water.
The upthrust is:

\mcoption{A}{$20\ \text{N}$}
\mcoption{B}{$8\ \text{N}$}
\mcoption{C}{$4\ \text{N}$}
\mcoption{D}{$12\ \text{N}$}
\ansref{Ch11-A3}

\vspace{0.4cm}

\textbf{A4.} A wooden block floats with $60\%$ of
its volume submerged. The density of the wood is
($\rho_{water} = 1000\ \text{kg\,m}^{-3}$):

\mcoption{A}{$400\ \text{kg\,m}^{-3}$}
\mcoption{B}{$600\ \text{kg\,m}^{-3}$}
\mcoption{C}{$1667\ \text{kg\,m}^{-3}$}
\mcoption{D}{$1000\ \text{kg\,m}^{-3}$}
\ansref{Ch11-A4}

\vspace{0.4cm}

\textbf{A5.} Pascal's principle is the basis for:

\mcoption{A}{Why ships float}
\mcoption{B}{How hydraulic lifts work}
\mcoption{C}{Why pressure increases with depth}
\mcoption{D}{How thermometers measure temperature}
\ansref{Ch11-A5}

\vspace{0.4cm}

\textbf{A6.} The relative density of a substance
is 8.5. Its density in kg\,m$^{-3}$ is:

\mcoption{A}{$8.5\ \text{kg\,m}^{-3}$}
\mcoption{B}{$850\ \text{kg\,m}^{-3}$}
\mcoption{C}{$8500\ \text{kg\,m}^{-3}$}
\mcoption{D}{$85\,000\ \text{kg\,m}^{-3}$}
\ansref{Ch11-A6}

\vspace{0.4cm}

\textbf{A7.} An object is neutrally buoyant when:

\mcoption{A}{It has zero mass}
\mcoption{B}{Its density equals the fluid density}
\mcoption{C}{The upthrust equals twice its weight}
\mcoption{D}{It is at the surface of the fluid}
\ansref{Ch11-A7}

\vspace{0.4cm}

\textbf{A8.} Water pressure at the bottom of two
containers --- one wide and one narrow --- filled
to the same depth is:

\mcoption{A}{Greater in the wider container}
\mcoption{B}{Greater in the narrower container}
\mcoption{C}{Equal in both containers}
\mcoption{D}{Depends on the total mass of water}
\ansref{Ch11-A8}

\vspace{0.4cm}

\textbf{A9.} A hydraulic press has pistons of area
$2\ \text{cm}^2$ and $200\ \text{cm}^2$. A force of
$10\ \text{N}$ on the small piston produces a force
on the large piston of:

\mcoption{A}{$0.1\ \text{N}$}
\mcoption{B}{$100\ \text{N}$}
\mcoption{C}{$1000\ \text{N}$}
\mcoption{D}{$10\,000\ \text{N}$}
\ansref{Ch11-A9}

\vspace{0.4cm}

\textbf{A10.} An iron anchor weighs $500\ \text{N}$
in air. When lowered into the sea, the tension in
the rope holding it is less than $500\ \text{N}$
because:

\mcoption{A}{The rope stretches and reduces the weight}
\mcoption{B}{The anchor dissolves slightly in sea water}
\mcoption{C}{An upthrust equal to the weight of sea
water displaced acts on the anchor}
\mcoption{D}{Water pressure reduces the effect of gravity}
\ansref{Ch11-A10}

\vspace{0.6cm}

\subsection*{Section B: Theory and Calculations}

\textbf{B1.} A rectangular barge has dimensions
$30\ \text{m} \times 10\ \text{m} \times 3\ \text{m}$
(length $\times$ width $\times$ depth). It floats
in fresh water ($\rho = 1000\ \text{kg\,m}^{-3}$)
with $1.5\ \text{m}$ of its depth submerged.
($g = 10\ \text{m\,s}^{-2}$)\\[4pt]
(a) Calculate the volume of water displaced.\\
(b) Calculate the upthrust on the barge.\\
(c) Calculate the mass of the barge (including cargo).\\
(d) The barge moves to sea water
($\rho = 1025\ \text{kg\,m}^{-3}$) with the same mass.
By how much does the depth of submersion change?
\hfill\ansref{Ch11-B1}

\vspace{0.4cm}

\textbf{B2.} An alloy is made by mixing
$200\ \text{g}$ of metal A (density
$8000\ \text{kg\,m}^{-3}$) with $300\ \text{g}$
of metal B (density $6000\ \text{kg\,m}^{-3}$),
assuming volumes are additive.\\[4pt]
(a) Calculate the volume of each metal.\\
(b) Calculate the density of the alloy.\\
(c) If a sphere of this alloy has mass $130\ \text{g}$,
calculate its radius.\\
(d) Would this sphere float in mercury
($\rho_{Hg} = 13\,600\ \text{kg\,m}^{-3}$)? Justify.
\hfill\ansref{Ch11-B2}

\vspace{0.4cm}

\textbf{B3.} A hydraulic braking system in a car has
a master cylinder of diameter $2.0\ \text{cm}$ connected
to four wheel cylinders, each of diameter $4.0\ \text{cm}$.
The driver applies a force of $200\ \text{N}$ on the
brake pedal which is connected to the master cylinder
via a lever with mechanical advantage of $4.0$.\\[4pt]
(a) Calculate the force applied to the master cylinder piston.\\
(b) Calculate the pressure in the brake fluid.\\
(c) Calculate the force exerted at each wheel cylinder.\\
(d) Calculate the total braking force on all four wheels.\\
(e) If the car has mass $1200\ \text{kg}$ and is
travelling at $20\ \text{m\,s}^{-1}$, how far does
it travel before stopping? (Assume the braking force
is the net decelerating force.)
\hfill\ansref{Ch11-B3}

\vspace{0.4cm}

\textbf{B4.} A stone weighs $2.50\ \text{N}$ in air.
When submerged in liquid X it weighs $2.10\ \text{N}$
and when submerged in water it weighs $2.00\ \text{N}$.\\[4pt]
(a) Calculate the upthrust in each liquid.\\
(b) Calculate the volume of the stone.\\
(c) Calculate the density of liquid X.\\
(d) Calculate the relative density of the stone.\\
(e) The stone is placed in liquid X in a closed container
and the container is pressurised. Does the stone sink
deeper, rise, or remain at the same depth? Explain
using Pascal's principle and Archimedes' principle.
\hfill\ansref{Ch11-B4}

\vspace{0.4cm}

\textbf{B5.} An iceberg floats in sea water with
$90\%$ of its volume below the surface. The density
of sea water is $1025\ \text{kg\,m}^{-3}$.\\[4pt]
(a) Calculate the density of ice.\\
(b) An expedition team stands on top of the iceberg.
Their combined mass is $600\ \text{kg}$. How much
further does the iceberg sink? (Assume the iceberg
has a uniform horizontal cross-sectional area of
$5000\ \text{m}^2$.)\\
(c) A cubic block of ice of side $2.0\ \text{m}$
melts completely. Using your calculated density of
ice, find the mass of the ice block.\\
(d) When the ice melts, does the sea level rise,
fall, or remain the same? Justify mathematically.
\hfill\ansref{Ch11-B5}

\end{exercisebox}
